\section{Decay on the tree}
\label{sec:tree-decay}

We now apply Proposition~\ref{prop:cone-closure} to the observations in
successive generations.  The saving from just two branches will suffice;
the remaining branches need only satisfy subadditivity.

For the calculation, supply an ordering of the children at every vertex.
In the Galton--Watson model, one may construct this ordered tree by giving
each vertex an independent offspring count and numbering its children
from $1$ to that count.  Erasing the numbering gives the unlabelled tree
in the theorem.  Write $\widehat T$ for the ordered tree and
$\widehat T_{\le n}$ for its truncation at depth $n$.  Define
\[
  Y_n=\bigl(\widehat T_{\le n},(\sigma_v)_{v\in L_n}\bigr),
  \qquad p^{(n)}_i=\Pr(\sigma_\rho=i\mid Y_n).
\]
Every finite-depth truncation is finite almost surely.  Let $\mu_n$ be
the law of $p^{(n)}$, and put
\begin{equation}
  m_n=\int_{\Deltafour}A(p)\,\mathrm d\mu_n(p).
  \label{eq:tree-moment}
\end{equation}
In the random-tree case this integral averages over the tree as well as
the spins, without conditioning on survival.

There are two information comparisons to check.  First, conditional on
$\widehat T_{\le n}$, the part of $\widehat T$ below depth $n$ is
independent of the root and all spins used in $Y_n$.  This follows from
the independence of the tree construction and the channel randomness.
Consequently, supplying the entire ordered tree leaves the posterior
unchanged:
\begin{equation}
  \Pr\bigl(\sigma_\rho=i\mid
       \widehat T,(\sigma_v)_{v\in L_n}\bigr)=p^{(n)}_i
  \quad\text{almost surely}.
  \label{eq:whole-tree-posterior}
\end{equation}
Second, forgetting the ordering cannot increase expected total variation
from the prior.  Indeed, for observation sigma-fields
$\mathcal F\subseteq\mathcal G$, the posterior given $\mathcal F$ is
the conditional expectation of the posterior given $\mathcal G$;
convexity of the total-variation norm gives the assertion.  Combining
this observation with \eqref{eq:quadratic-tv} and Jensen's inequality,
the advantage defined in Theorem~\ref{thm:main} satisfies
\begin{equation}
  a_n(d,\lambda)\le\frac12\E\sqrt{A(p^{(n)})}
  \le\frac12\sqrt{m_n}.
  \label{eq:tree-tv-bound}
\end{equation}
It therefore remains to prove that $m_n$ tends to zero.

\begin{proposition}[Moment recursion]
\label{prop:tree-recursion}
In either tree family of Theorem~\ref{thm:main}, the posterior law
$\mu_n$ is symmetric and satisfies $\int H\,\mathrm d\mu_n\ge0$ for
every $n\ge0$.  Let $D$ denote the offspring count, so $D=d$ on the
regular tree and $D$ has law $\operatorname{Poisson}(d)$ on the
Galton--Watson tree.  Then
\begin{equation}
  m_{n+1}\le d\lambda^2m_n
     -\frac{2}{1000}\Pr(D\ge2)(\lambda^2m_n)^3.
  \label{eq:tree-recursion}
\end{equation}
\end{proposition}

\begin{proof}
At depth zero the root spin is revealed.  Thus $\mu_0$ assigns mass
$1/4$ to each coordinate vector, is symmetric, and satisfies
$\int H\,\mathrm d\mu_0=13/100>0$.

Suppose $\mu_n\in\mathcal C$.
Conditional on $D=k$, each child contributes an experiment obtained by
passing a depth-$n$ experiment through one broadcast edge.  The branches
are independent conditional on the root spin, and their individual
posterior laws are all the pushforward $(C_\lambda)_*\mu_n$.  For the
Galton--Watson tree this uses the independence of the child subtrees
from each other and from $D$; for the regular tree the subtrees are
deterministic copies of the same rooted tree.  By
Lemma~\ref{lem:channel} and Proposition~\ref{prop:cone-closure}, each
branch law belongs to $\mathcal C$ and has mean $A$ equal to
\[
  b=\lambda^2m_n.
\]
Combining the branches successively by Lemma~\ref{lem:bayes-product}
and Proposition~\ref{prop:cone-closure} preserves membership in
$\mathcal C$.  If $k=0$, the posterior is $\unif$ and has $H=A=0$.
Since $D$ is independent of the root and is observed, mixing these laws
over $D$ also preserves membership in $\mathcal C$.  This proves
the induction assertion for $\mu_{n+1}$.

For $k\ge2$, combine the first two branches before adding the others.
Their individual posterior laws lie in $\mathcal C$ and both have mean
$A$ equal to $b$.  The final assertion of
Proposition~\ref{prop:cone-closure} therefore bounds the mean of $A$
after their combination by $2b-2b^3/1000$.  Each additional branch
increases this upper bound by at most $b$, since the combined experiment
still belongs to $\mathcal C$.  For $k=0$ the mean is zero, and for $k=1$ it is
$b$.  In all cases,
\[
  \E\bigl[A(p^{(n+1)})\mid D=k\bigr]
  \le kb-\frac{2b^3}{1000}\boldsymbol{1}_{\{k\ge2\}}.
\]
Taking expectation over $D$, whose mean is $d$, proves
\eqref{eq:tree-recursion}.
\end{proof}

\begin{proof}[Proof of Theorem~\ref{thm:main}]
If $\lambda=0$, each branch observation is independent of the root;
equivalently, \eqref{eq:tree-recursion} gives $m_n=0$ for every $n\ge1$.
Assume henceforth that $\lambda>0$ and $d\lambda^2\le1$.

On the regular tree, $d\ge2$ and hence $\Pr(D\ge2)=1$.  The recursion
gives
\[
  0\le m_{n+1}\le m_n-\frac{2\lambda^6}{1000}m_n^3.
\]
Thus $(m_n)$ is decreasing and has a limit $\ell\ge0$.  Passing to the
limit yields $\ell\le\ell-2\lambda^6\ell^3/1000$, so $\ell=0$.

For Poisson offspring,
\[
  \Pr(D\ge2)=1-(1+d)e^{-d}>0\qquad(d>0).
\]
The same argument applies with the positive coefficient
$2\lambda^6(1-(1+d)e^{-d})/1000$ in place of $2\lambda^6/1000$.
In particular, it covers $d\le1$, including the noiseless critical case
$d=\lambda=1$, directly in the unconditioned tree law.  No step requires
an infinite tree or a nonempty observed generation: when $L_n$ is empty,
the tree observation is independent of the root and the posterior is
$\unif$.

Both arguments allow $d\lambda^2=1$; the strictly positive cubic term
is what forces decay at that boundary.  Finally,
\eqref{eq:tree-tv-bound} gives $a_n(d,\lambda)\to0$ in each family.
\end{proof}
